%%%%%%%%%%%%%%%%%%%%%%%%%%%%%%%%%%%%%%%%%
% EUR Dissertation Template
% Based on kaobook with EUR Corporate Identity
%
% Compile with: LuaLaTeX or XeLaTeX (recommended)
%               pdfLaTeX also works but with limited font options
%
%%%%%%%%%%%%%%%%%%%%%%%%%%%%%%%%%%%%%%%%%

\documentclass[
    fontsize=11pt,
    a4paper,
    twoside=true,
    open=any,
    %secnumdepth=2,
]{kaobook}

%----------------------------------------------------------------------------------------
%	FONT CONFIGURATION
%	Change the option to switch fonts:
%	- libertinus (default, elegant serif)
%	- firasans (modern, technical)
%	- ibmplex (professional, readable)
%	- helvetica (classic sans-serif)
%	- museosans (if you have the font installed)
%----------------------------------------------------------------------------------------

% Font: Fira Sans (modern, excellent for technical/academic docs)
\usepackage[firasans]{kaofonts}

%----------------------------------------------------------------------------------------
%	EUR CORPORATE IDENTITY
%----------------------------------------------------------------------------------------

\usepackage{kaoci-eur}
\seteurlogopath{eur_ci/Erasmus University/}

%----------------------------------------------------------------------------------------
%	BIBLIOGRAPHY
%----------------------------------------------------------------------------------------

\usepackage{kaobiblio}
\usepackage{kaorefs}
\addbibresource{references.bib}

%----------------------------------------------------------------------------------------
%	DOCUMENT INFO
%----------------------------------------------------------------------------------------

\title{Navigating Sustainability Transitions}
\author{Your Name}
\date{\today}

% Dissertation-specific info
\promotor{Prof. dr. First Promotor\\Erasmus University Rotterdam}
\copromotor{Dr. Second Promotor\\Erasmus University Rotterdam}
\committee{%
Prof. dr. Committee Member 1, University Name\\
Prof. dr. Committee Member 2, University Name\\
Dr. Committee Member 3, University Name%
}
\defensedate{January 15, 2025}
\defensetime{13:00}
\defenselocation{Senaatszaal, Erasmus Building}

%----------------------------------------------------------------------------------------
%	BEGIN DOCUMENT
%----------------------------------------------------------------------------------------

\begin{document}

%----------------------------------------------------------------------------------------
%	TITLE PAGE
%----------------------------------------------------------------------------------------

% Option 1: Full EUR green title page
\eurtitlepage

% Option 2: Minimal title page (uncomment to use instead)
%\eurtitlepageminimal

%----------------------------------------------------------------------------------------
%	FRONT MATTER
%----------------------------------------------------------------------------------------

\frontmatter

% Committee page
\committeepage

% Table of contents
\tableofcontents

%----------------------------------------------------------------------------------------
%	MAIN MATTER
%----------------------------------------------------------------------------------------

\mainmatter
\setupeur % Apply EUR styling for chapters

%----------------------------------------------------------------------------------------
%	CHAPTER 1: INTRODUCTION
%----------------------------------------------------------------------------------------

\chapter{Introduction to Sustainability Transitions}
\labch{introduction}

\margintoc % Shows chapter contents in margin

Sustainability transitions represent fundamental shifts in socio-technical systems towards more sustainable modes of production and consumption.\sidenote{The multi-level perspective (MLP) provides a useful analytical framework for understanding these transitions.}

\section{The Multi-Level Perspective}

The MLP conceptualizes transitions as the interplay between three analytical levels:

\begin{enumerate}
    \item \textbf{Niches}: Protected spaces where radical innovations emerge
    \item \textbf{Regimes}: The dominant configuration of actors, institutions, and technologies
    \item \textbf{Landscape}: Exogenous macro-level developments
\end{enumerate}

\begin{eurbox}[title=Key Concept]
Transitions occur when niche innovations break through to challenge and eventually transform the regime, often facilitated by landscape pressures.
\end{eurbox}

\section{Research Questions}

This dissertation addresses the following questions:\marginnote{These questions build on established transition theory while addressing empirical gaps.}

\begin{enumerate}
    \item How do niche innovations scale up to challenge incumbent regimes?
    \item What role do intermediary actors play in transition processes?
    \item How can policy interventions accelerate sustainability transitions?
\end{enumerate}

%----------------------------------------------------------------------------------------
%	CHAPTER 2: WIDE VISUALIZATIONS EXAMPLE
%----------------------------------------------------------------------------------------

\chapter{Wide Visualizations and Data}
\labch{visualizations}

\margintoc

This chapter demonstrates wide figures and tables that span the full page width, perfect for complex diagrams and data presentations.

\section{Wide Figures}

The following figure spans the full text width plus margins, ideal for multi-level perspective diagrams:

\begin{widefigure}
    \centering
    % Placeholder - replace with your diagram
    \begin{tikzpicture}
        % Landscape level
        \fill[EURgreen!20] (0,6) rectangle (14,8);
        \node at (7,7) {\textbf{Landscape (macro-level)}};

        % Regime level
        \fill[EURgreen!40] (0,2.5) rectangle (14,5.5);
        \node at (7,4) {\textbf{Regime (meso-level)}};

        % Niche level
        \fill[EURgreen!60] (0,0) rectangle (14,2);
        \node at (7,1) {\textbf{Niches (micro-level)}};

        % Arrows showing dynamics
        \draw[->, thick, EURcyan] (3,1.5) -- (5,3) -- (7,5) -- (9,7);
        \node[right] at (9,7) {\small Transition pathway};
    \end{tikzpicture}
    \caption{The Multi-Level Perspective framework showing interactions between landscape, regime, and niche levels over time.}
    \labfig{mlp}
\end{widefigure}

\section{Wide Tables}

For comprehensive data presentation:\marginnote{Wide tables are essential for comparing multiple cases or presenting detailed quantitative analysis.}

\begin{widetable}
    \centering
    \caption{Comparison of transition pathways across different sustainability domains.}
    \labtab{pathways}
    \begin{tabular}{llllll}
        \toprule
        \textbf{Domain} & \textbf{Niche Innovation} & \textbf{Regime Response} & \textbf{Landscape Pressure} & \textbf{Timeline} & \textbf{Status} \\
        \midrule
        Energy & Solar PV & Gradual adoption & Climate policy & 30+ years & Ongoing \\
        Mobility & Electric vehicles & Resistance/co-opt & Air quality & 20+ years & Accelerating \\
        Food & Plant-based proteins & Diversification & Health concerns & 10+ years & Emerging \\
        Construction & Circular materials & Limited uptake & Resource scarcity & 5+ years & Early stage \\
        \bottomrule
    \end{tabular}
\end{widetable}

\section{Margin Figures}

For supporting visualizations, margin figures keep the main text flowing:

\begin{marginfigure}
    \centering
    \begin{tikzpicture}[scale=0.8]
        \draw[thick, EURgreen] (0,0) circle (1.5);
        \node at (0,0) {\small Niche};
        \draw[->, thick, EURcyan] (0,1.5) -- (0,2.5);
    \end{tikzpicture}
    \caption{Niche development cycle}
    \labfig{niche}
\end{marginfigure}

Regular text continues here with the margin figure alongside. This layout is particularly effective for presenting supplementary diagrams, data visualizations, or conceptual illustrations without interrupting the main narrative flow.

%----------------------------------------------------------------------------------------
%	CHAPTER 3: METHODOLOGY
%----------------------------------------------------------------------------------------

\chapter{Methodology}
\labch{methodology}

\margintoc

\section{Research Design}

This dissertation employs a mixed-methods research design combining:

\begin{eurbox}[title=Methods Overview]
\begin{itemize}
    \item Qualitative case studies of transition processes
    \item Quantitative network analysis of actor relationships
    \item Document analysis of policy frameworks
\end{itemize}
\end{eurbox}

\section{Case Selection}

Cases were selected based on:\sidenote{Case selection follows theoretical sampling principles as outlined by Eisenhardt (1989).}

\begin{enumerate}
    \item Variation in transition phase (early, accelerating, stabilizing)
    \item Geographic diversity (Netherlands, Germany, UK)
    \item Sectoral coverage (energy, mobility, food systems)
\end{enumerate}

\begin{eurboxgold}[title=Note on Data Collection]
All interviews were conducted between 2022-2024 with informed consent from participants. Interview protocols were approved by the EUR Ethics Review Board.
\end{eurboxgold}

%----------------------------------------------------------------------------------------
%	BACK MATTER
%----------------------------------------------------------------------------------------

\backmatter

\chapter*{Summary}
\addcontentsline{toc}{chapter}{Summary}

This dissertation examines sustainability transitions through the lens of the multi-level perspective, contributing both theoretical refinements and empirical insights from multiple case studies.

\chapter*{Samenvatting}
\addcontentsline{toc}{chapter}{Samenvatting}

Dit proefschrift onderzoekt duurzaamheidstransities vanuit het perspectief van de multi-level perspective, met zowel theoretische verfijningen als empirische inzichten uit meerdere case studies.

\chapter*{Acknowledgments}
\addcontentsline{toc}{chapter}{Acknowledgments}

I would like to thank my supervisors, colleagues, and family for their support throughout this journey.

%----------------------------------------------------------------------------------------
%	BIBLIOGRAPHY
%----------------------------------------------------------------------------------------

%\printbibliography[heading=bibintoc, title=References]

\end{document}
